\color{unidarkgreen}{\subsection{Автоматическое повторное включение (АПВ)}}
\color{black}

\begin{enumerate}

\item
Автоматическое повторное включение (АПВ) предназначено для быстрого восстановления первоначального состояния электрической сети после аварийного отключения путём автоматического включения выключателя.

\item
В случае присоединения линии к системе через два выключателя для каждого из них можно задействовать свою логику АПВ. Далее приведено описание АПВ для выключателя В1, логика АПВ для второго выключателя аналогична.

\item
Выбор действующих режимов АПВ для выключателя В1 реализуется тремя входными сигналами ФСУ:
\begin{itemize}
    \item \textit{«В1 ОНЛ + ННЭ1»} (включение с контролем отсутствия напряжения на линии и наличия напряжения на энергообъекте Э1) -- АПВ линии;
    \item \textit{«В1 ОНЭ1 + ННЛ»} (включение с контролем отсутствия напряжения на энергообъекте Э1 и наличия напряжения на линии) -- АПВ энергообъекта Э1 (шин);
    \item \textit{«В1 без контролей»} (включение без контроля напряжений).
\end{itemize}

\item
Возможные режимы АПВ приведены в таблице \ref{desc_apv:tbl1}.

\small
\setlength{\extrarowheight}{0.05cm} % высота ячеек в таблице
\begin{longtable}{|>{\centering\arraybackslash}m{4.9cm}|>{\centering\arraybackslash}m{3.7cm}|>{\centering\arraybackslash}m{3.7cm}|>{\centering\arraybackslash}m{3.7cm}|}
  \caption{Выбор режима АПВ\hfill\vspace{-0.5\baselineskip}}\label{desc_apv:tbl1}\\
  \hhline{|-|-|-|-|}
  \rowcolor{gray!30}
  & \multicolumn{3}{c|}{Входные сигналы ФСУ} \\
  \hhline{|-|-|-|-|}
  \rowcolor{gray!30}
  \textit{Режим АПВ} & \textit{АПВ В1 ОНЛ+ННЭ1} & \textit{АПВ В1 ОНЭ1+ННЛ} & \textit{АПВ В1 без контролей} \\
  \hhline{|-|-|-|-|}

\endfirsthead

\caption*{\hspace{3pt}\emph{Продолжение таблицы \ref{desc_apv:tbl1}\hfill\vspace{-0.5\baselineskip}}} \\ % сделано по ГОСТ 2.105 п.6.8.7
\hhline{|-|-|-|-|}
\rowcolor{gray!30}
  & \multicolumn{3}{c|}{Входные сигналы ФСУ} \\
\hhline{|-|-|-|-|}
\endhead
\hhline{|-|-|-|-|}

\endfoot
\endlastfoot

КС(УС) 			  			 & 0 & 0 & 0 \\  \hline
ОНЛ+ННЭ1 / КС(УС) 			 & 1 & 0 & 0 \\  \hline
ОНЭ1+ННЛ / КС(УС) 			 & 0 & 1 & 0 \\  \hline
ОНЛ+ННЭ1 / ОНЭ1+ННЛ / КС(УС) & 1 & 1 & 0 \\  \hline
Без контролей 				 & 0/1 & 0/1 & 1 \\  \hline

\end{longtable}
\setlength{\extrarowheight}{0.0cm} % возврат параметра высота ячеек в таблице
\normalsize

\item
Как видно из таблицы \ref{desc_apv:tbl1}, режим АПВ с КС(УС) введен всегда и выводится посредством входного сигнала ФСУ \textit{«В1 без контролей»}. А режимы ОНЛ+ННЭ1 и ОНЭ1+ННЛ могут быть введены одновременно. Режим без контролей применяется только для АПВ линии. Условия включения в различных режимах более подробно описаны в пункте \ref{func:ksn} настоящего руководства по эксплуатации.

\item
Пуск АПВ линии осуществляется по цепи несоответствия при аварийном отключении выключателя (по наличию входного сигнала ФСУ \textit{«Сигнал несоотв.»}). Срабатывание происходит при набранной выдержке времени готовности АПВ с контролем заданного режима и выполнения условий включения. Выдержка времени готовности АПВ после оперативного включения и первого цикла регулируется уставкой \textbf{«Тгот АПВ В1»}, после второго цикла -- \textbf{«Твосст гот. АПВ В1»}.

\item
При приеме сигнала о срабатывании ДЗШ (входной сигнал ФСУ \textit{«Сброс АПВл от ДЗШ»}) готовность АПВ линии сбрасывается.

\item
АПВ линии может быть выполнено двукратным, ввод второго цикла осуществляется программным ключом \textbf{«2-ой цикл АПВл В1»}. Первый цикл осуществляется при выполнении условий включения с выдержкой времени \textbf{«Тср АПВл-1 В1»}. В случае сброса сигнала пуска АПВ при отсчете времени \textbf{«Тср АПВл-1 В1»} цепь логики сбрасывается, и схема АПВ возвращается в исходное состояние (готовность АПВ с набранной выдержкой готовности). Если первый цикл АПВ был успешным, то начинается набор выдержки времени готовности к повторному действию.

\item
Если после первого цикла АПВ за время \textbf{«Тгот АПВ В1»} (выдержка времени к повторному действию не успела набраться) повторно поступает сигнал пуска АПВ, то осуществляется 2-ой цикл с выдержкой времени \textbf{«Тср АПВл-2 В1»}. В случае сброса сигнала пуска АПВ при отсчете времени \textbf{«Тср АПВл-2 В1»} цепь логики сбрасывается, и схема возвращается в режим набора выдержки времени готовности к повторному действию. Если второй цикл АПВ был успешным, то начинается набор выдержки времени готовности к повторному действию.

\item
Второй цикл АПВ может быть выведен от внешнего сигнала \textit{«Вывод АПВл-2 В1»}.

\item
Для схем с двумя выключателями на присоединение предусмотрен входной сигнал ФСУ \textit{«В1 ведомый»} для оперативного изменения порядка включения выключателя от АПВ линии – \textit{«ведущий»} или \textit{«ведомый»}. По умолчанию выключатель принимается “ведущим”. При наличии сигнала \textit{«В1 ведомый»} время включения от АПВ линии дополнительно увеличивается на выдержку времени \textbf{«Тср АПВл В1 доп.»}.

\item
Длительность импульса включения выключателя В1 от АПВ линии задается уставкой \textbf{«Тимп АПВл В1»}.

\item
Функциональная схема логики АПВ В1 ЛЭП приведена на рисунке \ref{pict_apv:APV_LAP}. Функциональная схема логики АПВ В2 ЛЭП выполняется аналогично схеме логики АПВ В1 ЛЭП.

\vspace{2mm}
\begin{figure}[H]
\centering
\includegraphics[width=0.9\textwidth,height=0.9\textheight,keepaspectratio]{\fbpath/03.01 ЛВ/761.3011 АПВ/img/APV_LAP.pdf}
\captionof{figure}{Функциональная схема логики АПВ В1 ЛЭП}\label{pict_apv:APV_LAP}
\end{figure}

\item
Уставки АПВ В1 ЛЭП представлены в таблице \ref{app_set:apv_q1}. Уставки АПВ В2 ЛЭП аналогичны уставкам АПВ В1 ЛЭП.

\item
Пуск АПВ шин осуществляется от внешнего сигнала о срабатывании ДЗШ (входной сигнал ФСУ \textit{«Пуск АПВ В1 от ДЗШ»}). При этом АПВ линии запрещается. АПВ шин выполнено однократным и осуществляется при выполнении условий включения с выдержкой времени \textbf{«Тср АПВш В1»}. Длительность импульса включения выключателя В1 от АПВ шин задается уставкой \textbf{«Тимп АПВш В1»}.

\item
Функциональная схема логики АПВ В1 СШ приведена на рисунке \ref{pict_apv:APV_CSH}. Функциональная схема логики АПВ В2 СШ выполняется аналогично схеме логики АПВ В1 СШ. Уставки АПВ В1 СШ представлены в таблице \ref{app_set:apv_q1}. Уставки АПВ В2 СШ аналогичны уставкам АПВ В1 СШ.

\vspace{2mm}
\begin{figure}[H]
\centering
\includegraphics[width=0.9\textwidth,height=0.9\textheight,keepaspectratio]{\fbpath/03.01 ЛВ/761.3011 АПВ/img/APV_CSH.pdf}
\captionof{figure}{Функциональная схема АПВ В1 СШ}\label{pict_apv:APV_CSH}
\end{figure}

\item
В устройстве предусмотрено ограничение ожидания условий включения. Если в течение выдержки времени \textbf{«Тожид усл. вкл. В1»} после аварийного отключения не возникают условия включения, то действие АПВ запрещается.

\item
Предусмотрена возможность запрета АПВ при длительно отключенном выключателе, который вводится программным ключом \textbf{«Огр. ожид. пуска АПВ В1»}. Действие на запрет АПВ осуществляется с выдержкой времени \textbf{«Тожид усл. пуска АПВ В1»}.

\item
Кроме того, выставлением соответствующих уставок может осуществляться запрет АПВ при срабатывании основной защиты, ступеней ДЗ-ФЗ и ДЗ-ФФ, ступеней ТЗНП, ступеней АМТЗ, при 
срабатывании оперативного ускорения (как с выдержкой, так и без выдержки времени) ДЗ и ТЗНП, поперечного ускорения ТЗНП, автоматического ускорения, ЗНР, МФТО, ТЗО, телеускорения и телеотключения, ЗНФ. На запрет АПВ также жёстко действуют срабатывание внутреннего или внешнего УРОВ, оперативное отключение, неуспешное АПВ смежного выключателя (если смежный выключатель «ведущий»), перевод режима управления выключателем в \textit{«Местное»}, отключение от аварийного снижения давления элегаза в ТТ. Также для запрета АПВ предусмотрен входной сигнал ФСУ \textit{«Внеш. запрет АПВ»}.

Логика работы запрета АПВ приведена на рисунках \ref{pict_apv:APV_ZAP_Q1}, \ref{pict_apv:APV_ZAP}.


\vspace{2mm}
\begin{figure}[H]
\centering
\includegraphics[width=0.9\textwidth,height=0.9\textheight,keepaspectratio]{\fbpath/03.01 ЛВ/761.3011 АПВ/img/APV_ZAP_Q1.pdf}
\captionof{figure}{Функциональная схема запрета АПВ В1}\label{pict_apv:APV_ZAP_Q1}
\end{figure}

\vspace{2mm}
\begin{figure}[H]
\centering
\includegraphics[width=0.9\textwidth,height=0.9\textheight,keepaspectratio]{\fbpath/03.01 ЛВ/761.3011 АПВ/img/APV_ZAP.pdf}
\captionof{figure}{Функциональная схема общего запрета АПВ}\label{pict_apv:APV_ZAP}
\end{figure}


\end{enumerate}